\section{第 2.1 节：最小作用量原理与拉格朗日力学}
\label{section:QM-C02-S01}

\studymeta
  {QM-C02-S01}
  {2026-08-27}
  {Shankar, \textit{Principles of Quantum Mechanics}, Second Edition}
  {第 2 章 \S 2.1；印刷页 pp.~75--83；PDF pp.~88--96}
  {内容、教材例题与练习 2.1.1--2.1.3 已核对}

\subsection{本节主线}

Newton 力学从某一时刻的位置与速度出发，用局部微分方程逐步确定轨迹；Lagrange 力学则给连接固定时空端点的每条候选路径分配作用量，并以作用量的一阶变分为零选出经典路径。变分计算产生 Euler--Lagrange 方程，而其形式不依赖所选的广义坐标。正则动量、循环坐标和守恒量也由这一结构自然出现。

\notetopic{QM-C02-S01-T01}{局部运动方程与全局路径原理}

一维粒子在与速度无关的势能 $V(x)$ 中满足 Newton 方程
\begin{equation}
  m\ddot x(t)=-\frac{dV}{dx}.
  \label{eq:QM-C02-S01-newton}
\end{equation}
给定初值 $x(t_i)$ 与 $\dot x(t_i)$ 后，式~\eqref{eq:QM-C02-S01-newton} 逐时刻确定后续运动，因此这是 local initial-value formulation（局部初值表述）。对 $n$ 个 Cartesian coordinates，系统的位形由
\[
  (x_1,\ldots,x_n)
\]
描述，它是 $n$ 维 configuration space（位形空间）中的一个点；系统的演化则是位形空间中的一条路径。

Lagrange 表述改为固定两个时空端点
\begin{equation}
  x(t_i)=x_i,
  \qquad
  x(t_f)=x_f,
  \label{eq:QM-C02-S01-fixed-endpoints}
\end{equation}
并比较所有连接它们的足够光滑路径。对速度无关的保守势，定义 Lagrangian（拉格朗日量）
\begin{equation}
  \mathcal L(x,\dot x,t)=T-V
  =\frac12m\dot x^2-V(x,t),
  \label{eq:QM-C02-S01-lagrangian}
\end{equation}
以及 action（作用量）
\begin{equation}
  \boxed{
  S[x]
  =\int_{t_i}^{t_f}\mathcal L(x(t),\dot x(t),t)\,dt.}
  \label{eq:QM-C02-S01-action}
\end{equation}
$\mathcal L$ 的输入是某一时刻的若干数，因而是普通函数；$S$ 的输入是整条路径函数 $x(t)$，因而是 functional（泛函）。作用量的量纲为能量乘时间，与角动量相同。

\begin{figure}[htbp]
  \centering
  \begin{tikzpicture}[x=1.15cm,y=0.9cm,>=Latex]
    \draw[->] (0,0)--(6.3,0) node[right] {$t$};
    \draw[->] (0,0)--(0,4.2) node[above] {$x$};
    \coordinate (A) at (0.7,0.8);
    \coordinate (B) at (5.6,3.4);
    \fill (A) circle (2pt) node[below right] {$(t_i,x_i)$};
    \fill (B) circle (2pt) node[above left] {$(t_f,x_f)$};
    \draw[noteBlue,very thick]
      plot[smooth] coordinates {(0.7,0.8) (1.8,1.3) (3.0,2.5) (4.2,2.6) (5.6,3.4)};
    \draw[verifyOrange,very thick,dashed]
      plot[smooth] coordinates {(0.7,0.8) (1.8,2.0) (3.0,1.7) (4.2,3.3) (5.6,3.4)};
    \node[noteBlue] at (3.9,2.25) {$x_{\mathrm{cl}}(t)$};
    \node[verifyOrange] at (2.4,3.15) {$x_{\mathrm{cl}}(t)+\varepsilon\eta(t)$};
  \end{tikzpicture}
  \caption{固定同一对端点的经典路径与一条允许的邻近路径。}
  \label{fig:QM-C02-S01-path-variation}
\end{figure}

经典路径 $x_{\mathrm{cl}}(t)$ 满足 stationary-action condition（平稳作用量条件）
\begin{equation}
  \boxed{\delta S[x_{\mathrm{cl}}]=0.}
  \label{eq:QM-C02-S01-stationary-action}
\end{equation}
它表示任意允许的无穷小路径变形都不产生作用量的一阶变化，而不保证作用量一定是全局最小值。传统名称 principle of least action（最小作用量原理）因此不如 principle of stationary action（平稳作用量原理）准确。

自由粒子给出一个确实取到最小值的检验。令 $V=0$，连接固定端点的匀速路径为
\[
  x_{\mathrm{cl}}(t)
  =x_i+\frac{x_f-x_i}{t_f-t_i}(t-t_i).
\]
对任意满足 $\eta(t_i)=\eta(t_f)=0$ 的邻近路径 $x=x_{\mathrm{cl}}+\eta$，有
\begin{align}
  S[x]-S[x_{\mathrm{cl}}]
  &=m\int_{t_i}^{t_f}\dot x_{\mathrm{cl}}\dot\eta\,dt
    +\frac m2\int_{t_i}^{t_f}\dot\eta^2\,dt \notag\\
  &=m\dot x_{\mathrm{cl}}[\eta(t_f)-\eta(t_i)]
    +\frac m2\int_{t_i}^{t_f}\dot\eta^2\,dt \notag\\
  &=\frac m2\int_{t_i}^{t_f}\dot\eta^2\,dt\geq0.
  \label{eq:QM-C02-S01-free-particle-minimum}
\end{align}
这个例子中的匀速直线是最小作用量路径；一般问题只能由一阶变分推出平稳性。

\notetopic{QM-C02-S01-T02}{Euler--Lagrange 方程的完整变分推导}

在固定端点的经典路径附近构造一族路径
\begin{equation}
  x_\varepsilon(t)=x_{\mathrm{cl}}(t)+\varepsilon\eta(t),
  \qquad
  \eta(t_i)=\eta(t_f)=0.
  \label{eq:QM-C02-S01-varied-path}
\end{equation}
于是
\[
  \dot x_\varepsilon(t)
  =\dot x_{\mathrm{cl}}(t)+\varepsilon\dot\eta(t).
\]
把泛函限制在这族路径上，$S[x_\varepsilon]$ 成为普通参数 $\varepsilon$ 的函数。平稳条件为
\begin{equation}
  \left.\frac{d}{d\varepsilon}S[x_\varepsilon]
  \right|_{\varepsilon=0}=0.
  \label{eq:QM-C02-S01-epsilon-stationary}
\end{equation}
对式~\eqref{eq:QM-C02-S01-action} 使用链式法则，得到一阶变分
\begin{align}
  \delta S
  &=\left.\frac{d}{d\varepsilon}
    \int_{t_i}^{t_f}
    \mathcal L(x_\varepsilon,\dot x_\varepsilon,t)\,dt
    \right|_{\varepsilon=0} \notag\\
  &=\int_{t_i}^{t_f}
    \left(
      \frac{\partial\mathcal L}{\partial x}\eta
      +\frac{\partial\mathcal L}{\partial\dot x}\dot\eta
    \right)dt.
  \label{eq:QM-C02-S01-first-variation}
\end{align}
第二项含 $\dot\eta$，对其分部积分：
\begin{align}
  \int_{t_i}^{t_f}
  \frac{\partial\mathcal L}{\partial\dot x}\dot\eta\,dt
  &=\left[
    \frac{\partial\mathcal L}{\partial\dot x}\eta
    \right]_{t_i}^{t_f}
  -\int_{t_i}^{t_f}
    \frac{d}{dt}
    \left(\frac{\partial\mathcal L}{\partial\dot x}\right)
    \eta\,dt.
  \label{eq:QM-C02-S01-variation-integration-by-parts}
\end{align}
固定端点使边界项严格为零，故
\begin{equation}
  \delta S
  =\int_{t_i}^{t_f}
  \left[
    \frac{\partial\mathcal L}{\partial x}
    -\frac{d}{dt}
    \left(\frac{\partial\mathcal L}{\partial\dot x}\right)
  \right]\eta(t)\,dt.
  \label{eq:QM-C02-S01-variation-after-parts}
\end{equation}
若上式对所有内部任意、端点为零的 $\eta(t)$ 都等于零，则变分法基本引理要求方括号内逐点为零。否则若该函数在某个小邻域内同号，可选择只支撑在该邻域且同号的 $\eta$，使积分不为零。由此得到 Euler--Lagrange equation：
\begin{equation}
  \boxed{
  \frac{d}{dt}
  \left(\frac{\partial\mathcal L}{\partial\dot x}\right)
  -\frac{\partial\mathcal L}{\partial x}=0.}
  \label{eq:QM-C02-S01-euler-lagrange}
\end{equation}
逻辑上必须要求 $\delta S=0$ 对所有允许变分成立；单独某个 $\eta$ 使积分为零，可能只是正负贡献抵消，不能推出被积函数为零。

对 $n$ 个独立坐标，逐个变分可得
\begin{equation}
  \boxed{
  \frac{d}{dt}
  \left(\frac{\partial\mathcal L}{\partial\dot q_i}\right)
  -\frac{\partial\mathcal L}{\partial q_i}=0,
  \qquad i=1,\ldots,n.}
  \label{eq:QM-C02-S01-multicoordinate-euler-lagrange}
\end{equation}
推导只使用独立坐标、链式法则、分部积分和固定端点条件，并未要求 $q_i$ 是 Cartesian coordinates。

把式~\eqref{eq:QM-C02-S01-lagrangian} 代入式~\eqref{eq:QM-C02-S01-euler-lagrange}，有
\[
  \frac{\partial\mathcal L}{\partial\dot x}=m\dot x,
  \qquad
  \frac{\partial\mathcal L}{\partial x}=-\frac{dV}{dx},
\]
所以
\[
  m\ddot x=-\frac{dV}{dx},
\]
恢复 Newton 第二定律。简谐振子的直接应用见\relatedexercise{QM-C02-S01-E01}{教材练习 2.1.1}。

\notetopic{QM-C02-S01-T03}{广义坐标、正则动量与循环坐标}

能够无冗余地描述系统位形的独立参数 $q_i$ 称为 generalized coordinates（广义坐标）。它们可以是长度、角度或其他受约束系统的自然参数。由于式~\eqref{eq:QM-C02-S01-multicoordinate-euler-lagrange} 在任意这类坐标中形式不变，实际计算时可以选择最能反映几何约束和对称性的坐标。

定义与 $q_i$ 共轭的 canonical momentum（正则动量）
\begin{equation}
  \boxed{p_i=\frac{\partial\mathcal L}{\partial\dot q_i}}
  \label{eq:QM-C02-S01-canonical-momentum}
\end{equation}
后，Euler--Lagrange 方程可简写为
\begin{equation}
  \dot p_i=\frac{\partial\mathcal L}{\partial q_i}.
  \label{eq:QM-C02-S01-generalized-newton}
\end{equation}
右侧一般不能直接解释为物理力。标准的 generalized applied force（广义作用力）$Q_i^{\mathrm{appl}}$ 满足
\begin{equation}
  \boxed{
  \frac{d}{dt}\left(\frac{\partial T}{\partial\dot q_i}\right)
  -\frac{\partial T}{\partial q_i}
  =Q_i^{\mathrm{appl}}.}
  \label{eq:QM-C02-S01-generalized-force}
\end{equation}
若作用力来自速度无关的普通势能，则 $Q_i^{\mathrm{appl}}=-\partial V/\partial q_i$，式~\eqref{eq:QM-C02-S01-generalized-force} 与 $\mathcal L=T-V$ 的 Euler--Lagrange 方程等价。在这一保守势情形下有
\[
  \frac{\partial\mathcal L}{\partial q_i}
  =\frac{\partial T}{\partial q_i}+Q_i^{\mathrm{appl}},
\]
所以 $\partial\mathcal L/\partial q_i$ 还可能包含曲线坐标在动能中产生的几何项。只有在 Cartesian coordinates 且 $T$ 不显含位置时，它才直接退化为相应的物理力分量。$p_i$ 也不一定是 $m\dot q_i$；例如若 $q_i$ 是角度，对应的 $p_i$ 可为角动量。正则动量必须由式~\eqref{eq:QM-C02-S01-canonical-momentum} 计算，不能从名称猜测。

若完整 Lagrangian 不显含某个坐标 $q_j$，即
\begin{equation}
  \frac{\partial\mathcal L}{\partial q_j}=0,
  \label{eq:QM-C02-S01-cyclic-condition}
\end{equation}
则 $q_j$ 称为 cyclic coordinate（循环坐标），并有
\begin{equation}
  \boxed{p_j=\text{constant}.}
  \label{eq:QM-C02-S01-cyclic-conservation}
\end{equation}
判断时必须检查完整的 $\mathcal L$，不能只检查势能。坐标可以不出现在 $\mathcal L$ 中，而其速度 $\dot q_j$ 仍然出现；这正是相应正则动量非零且守恒的典型情形。

\notetopic{QM-C02-S01-T04}{教材例题 2.1.1：平面极坐标}

令
\[
  x=\rho\cos\phi,
  \qquad
  y=\rho\sin\phi.
\]
微分后平方相加：
\begin{align}
  dx&=\cos\phi\,d\rho-\rho\sin\phi\,d\phi,\notag\\
  dy&=\sin\phi\,d\rho+\rho\cos\phi\,d\phi,\notag\\
  ds^2=dx^2+dy^2&=d\rho^2+\rho^2d\phi^2.
  \label{eq:QM-C02-S01-polar-line-element}
\end{align}
因此速度平方和 Lagrangian 为
\begin{equation}
  v^2=\dot\rho^2+\rho^2\dot\phi^2,
  \qquad
  \mathcal L
  =\frac12m(\dot\rho^2+\rho^2\dot\phi^2)-V(\rho,\phi).
  \label{eq:QM-C02-S01-polar-lagrangian}
\end{equation}

\begin{figure}[htbp]
  \centering
  \begin{tikzpicture}[scale=0.95,>=Latex]
    \draw[->] (-0.2,0)--(5.4,0) node[right] {$x$};
    \draw[->] (0,-0.2)--(0,4.2) node[above] {$y$};
    \coordinate (P) at (4,2.7);
    \draw[noteBlue,very thick,->] (0,0)--(P) node[midway,below right] {$\rho$};
    \draw[->] (1.15,0) arc[start angle=0,end angle=34,radius=1.15];
    \node at (1.35,0.45) {$\phi$};
    \fill (P) circle (2pt);
    \draw[->,verifyOrange,thick] (P)--($(P)+(1.2,0.81)$) node[right] {$\hat{\bm e}_\rho$};
    \draw[->,noteBlue,thick] (P)--($(P)+(-0.81,1.2)$) node[above] {$\hat{\bm e}_\phi$};
    \draw[dashed] (P)--(4,0);
  \end{tikzpicture}
  \caption{极坐标的径向与角向基矢；基矢随 $\phi$ 改变。}
  \label{fig:QM-C02-S01-polar-basis}
\end{figure}

径向坐标给出
\begin{align}
  \frac{\partial\mathcal L}{\partial\dot\rho}&=m\dot\rho,
  &
  \frac{\partial\mathcal L}{\partial\rho}
  &=m\rho\dot\phi^2-\frac{\partial V}{\partial\rho},
\end{align}
故
\begin{equation}
  \boxed{
  m(\ddot\rho-\rho\dot\phi^2)
  =-\frac{\partial V}{\partial\rho}.}
  \label{eq:QM-C02-S01-polar-radial}
\end{equation}
角坐标给出
\begin{equation}
  \boxed{
  \frac{d}{dt}(m\rho^2\dot\phi)
  =-\frac{\partial V}{\partial\phi},}
  \label{eq:QM-C02-S01-polar-angular-conservative}
\end{equation}
或展开为
\begin{equation}
  \boxed{
  m(\rho\ddot\phi+2\dot\rho\dot\phi)
  =-\frac1\rho\frac{\partial V}{\partial\phi}.}
  \label{eq:QM-C02-S01-polar-angular}
\end{equation}
$-\rho\dot\phi^2$ 和 $2\dot\rho\dot\phi$ 来自曲线坐标基的变化；Lagrange 方法把这些几何信息编码在动能中，无须单独对旋转基矢求导。
特别地，$\partial\mathcal L/\partial\rho=m\rho\dot\phi^2-\partial V/\partial\rho$ 中，$m\rho\dot\phi^2$ 是动能的坐标导数，不是额外的径向作用力。真正的径向物理力为 $-\partial V/\partial\rho$，而完整径向加速度是 $\ddot\rho-\rho\dot\phi^2$。

若 $V=V(\rho)$，则 $\phi$ 是循环坐标，
\begin{equation}
  p_\phi=m\rho^2\dot\phi=L_z=\text{constant}.
  \label{eq:QM-C02-S01-polar-angular-momentum}
\end{equation}
相应的面积速度为
\begin{equation}
  \frac{dA}{dt}=\frac12\rho^2\dot\phi=\frac{L_z}{2m},
  \label{eq:QM-C02-S01-areal-velocity}
\end{equation}
所以中心力运动满足等时扫过等面积。教材例题还表明，适当选取坐标后，守恒量常可由“某坐标不显含于 $\mathcal L$”直接读出。

\notetopic{QM-C02-S01-T05}{三维中心势中的球坐标与角动量}

球坐标的线元为
\begin{equation}
  ds^2=dr^2+r^2d\theta^2+r^2\sin^2\theta\,d\phi^2.
  \label{eq:QM-C02-S01-spherical-line-element}
\end{equation}
对中心势 $V(r)$，
\begin{equation}
  \mathcal L
  =\frac12m\left(
    \dot r^2+r^2\dot\theta^2
    +r^2\sin^2\theta\,\dot\phi^2
  \right)-V(r).
  \label{eq:QM-C02-S01-spherical-lagrangian}
\end{equation}
虽然 $V$ 与 $\theta$ 无关，但动能显含 $\sin^2\theta$，所以 $\theta$ 不是循环坐标。$\phi$ 不显含于完整的 $\mathcal L$，因此
\begin{equation}
  p_\phi=mr^2\sin^2\theta\,\dot\phi=L_z
  \label{eq:QM-C02-S01-spherical-p-phi}
\end{equation}
守恒。球坐标选定的 $z$ 轴使绕该轴的旋转表现为 $\phi\mapsto\phi+\alpha$，所以 $L_z$ 的守恒最直接；中心势的完整旋转对称性实际上还保证整个角动量向量守恒：
\begin{equation}
  \frac{d\bm L}{dt}=\bm r\times\bm F=0,
  \qquad
  L^2=p_\theta^2+\frac{p_\phi^2}{\sin^2\theta}.
  \label{eq:QM-C02-S01-full-angular-momentum}
\end{equation}
球坐标下的三条运动方程及循环坐标判断见\relatedexercise{QM-C02-S01-E03}{教材练习 2.1.3}。

\subsection{适用条件与常见误区}

\begin{itemize}[leftmargin=*]
  \item $\delta S=0$ 只说明一阶变分为零，不自动说明作用量取得全局最小值。
  \item 固定端点条件用于消去分部积分产生的边界项；若端点允许变化，还会出现额外的自然边界条件。
  \item $\mathcal L=T-V$ 适用于本节的速度无关保守势。后续电磁相互作用表明，一般 Lagrangian 不必具有这一形式。
  \item 正则动量由 $p_i=\partial\mathcal L/\partial\dot q_i$ 定义，不一定等于机械动量 $m\dot q_i$。
  \item 判断循环坐标必须检查完整的 $\mathcal L$，不能只检查 $V$。
\end{itemize}

\subsection{一分钟回顾}

作用量是以整条路径为输入的泛函，经典路径使其一阶变分为零。固定端点和分部积分把平稳条件转化为 Euler--Lagrange 方程，其形式对任意合适的广义坐标保持不变。正则动量是 Lagrangian 对广义速度的偏导；若某坐标不显含于完整 Lagrangian，对应正则动量守恒。极坐标例题展示了曲线坐标的几何项如何由动能自动产生，也把旋转对称性与角动量守恒联系起来。

\section{第 2.2 节：电磁场中的 Lagrangian}
\label{section:QM-C02-S02}

\studymeta
  {QM-C02-S02}
  {2026-08-28}
  {Shankar, \textit{Principles of Quantum Mechanics}, Second Edition}
  {第 2 章 \S 2.2；印刷页 pp.~83--84；PDF pp.~96--97}
  {正文已核对；本节无教材例题或习题，配套题明确标为自编检验题}

\subsection{本节主线}

普通位置势 $V(\bm r,t)$ 只能产生梯度力，无法表示显式依赖速度的磁力。本节把 scalar potential（标势）$\phi$ 与 vector potential（矢势）$\bm A$ 写入 Lagrangian，并由 Euler--Lagrange 方程恢复完整 Lorentz force（洛伦兹力）。这一过程还迫使我们区分 canonical momentum（正则动量）与 mechanical momentum（机械动量），也说明“外力”“约束力”与“是否写在 Euler--Lagrange 方程右侧”不是同一分类标准。

\notetopic{QM-C02-S02-T01}{普通位置势的不足与电磁 Lagrangian}

教材使用 Gaussian units（高斯单位制），其中 Lorentz force 为
\begin{equation}
  \boxed{
  \bm F=q\left(\bm E+\frac{\bm v}{c}\times\bm B\right).}
  \label{eq:QM-C02-S02-lorentz-force}
\end{equation}
若 $\mathcal L=\frac12m\bm v^2-V(\bm r,t)$，则 Euler--Lagrange 方程只能给出 $m\dot{\bm v}=-\nabla V$，无法产生速度依赖项 $q\bm v\times\bm B/c$。正确的电磁 Lagrangian 是
\begin{equation}
  \boxed{
  \mathcal L_{\mathrm{em}}
  =\frac12m\bm v^2-q\phi(\bm r,t)
  +\frac qc\bm v\cdot\bm A(\bm r,t).}
  \label{eq:QM-C02-S02-em-lagrangian}
\end{equation}
$\bm v\cdot\bm A$ 是旋转下的标量；它对速度求偏导留下 $\bm A$，对位置求偏导产生 $\bm A$ 的空间导数，二者在 Euler--Lagrange 方程中组合成 $\nabla\times\bm A$。在 SI units 中，式~\eqref{eq:QM-C02-S02-lorentz-force} 和式~\eqref{eq:QM-C02-S02-em-lagrangian} 中相应的 $1/c$ 均不出现。

\notetopic{QM-C02-S02-T02}{标势、矢势与规范自由度}

势的表示不是从 Lagrangian 中倒推定义出来的，而是源于 homogeneous Maxwell equations（齐次 Maxwell 方程）
\begin{equation}
  \nabla\cdot\bm B=0,
  \qquad
  \nabla\times\bm E=-\frac1c\frac{\partial\bm B}{\partial t}.
  \label{eq:QM-C02-S02-homogeneous-maxwell}
\end{equation}
在拓扑足够简单的区域内，散度为零的光滑向量场可局部写成另一个向量场的旋度，因此
\begin{equation}
  \boxed{\bm B=\nabla\times\bm A.}
  \label{eq:QM-C02-S02-vector-potential}
\end{equation}
根据 Stokes theorem，$\oint_{\partial S}\bm A\cdot d\bm l=\int_S\bm B\cdot d\bm S$，所以 $\bm A$ 沿闭合曲线的环流编码了曲线内部的磁通量。
\begin{figure}[htbp]
  \centering
  \begin{tikzpicture}[scale=0.9,>=Latex]
    \fill[noteBlue!8] (0,0) ellipse (3.0 and 1.35);
    \draw[noteBlue,very thick] (0,0) ellipse (3.0 and 1.35);
    \draw[->,noteBlue,very thick]
      (2.82,-0.45) arc[start angle=-20,end angle=32,x radius=3.0,y radius=1.35];
    \node[noteBlue] at (3.55,0.55) {$C=\partial S$};
    \node at (-2.35,-0.8) {$S$};
    \foreach \x/\y in {-1.8/0.35,-0.9/-0.45,0/0.45,0.9/-0.35,1.8/0.3}
      {\node[verifyOrange] at (\x,\y) {$\odot$};}
    \node[verifyOrange] at (0,-1.85) {$\bm B$ 穿过曲面};
  \end{tikzpicture}
  \caption{Stokes theorem 把边界上的矢势环流与穿过曲面的磁通量联系起来；箭头只表示闭合曲线的正向。}
  \label{fig:QM-C02-S02-stokes}
\end{figure}
把式~\eqref{eq:QM-C02-S02-vector-potential} 代入 Faraday's law，有
\[
  \nabla\times\left(
    \bm E+\frac1c\frac{\partial\bm A}{\partial t}
  \right)=0.
\]
无旋场可局部写成标量势的负梯度，故
\begin{equation}
  \boxed{
  \bm E=-\nabla\phi-\frac1c\frac{\partial\bm A}{\partial t}.}
  \label{eq:QM-C02-S02-scalar-potential}
\end{equation}
当磁场随时间变化时，$\nabla\times\bm E$ 一般非零，因而不能只写 $\bm E=-\nabla\phi$。

势并不唯一。对任意足够光滑的 $\chi(\bm r,t)$，gauge transformation（规范变换）
\begin{equation}
  \bm A'=\bm A+\nabla\chi,
  \qquad
  \phi'=\phi-\frac1c\frac{\partial\chi}{\partial t}
  \label{eq:QM-C02-S02-gauge-transform}
\end{equation}
保持 $\bm E$ 与 $\bm B$ 不变，同时令
\begin{equation}
  \mathcal L_{\mathrm{em}}'
  =\mathcal L_{\mathrm{em}}+\frac qc\frac{d\chi}{dt}.
  \label{eq:QM-C02-S02-gauge-lagrangian}
\end{equation}
作用量只增加固定端点处的边界值 $\frac qc[\chi_f-\chi_i]$，所以 Euler--Lagrange 方程不变。这把规范自由度与第 2.1 节的固定端点变分直接联系起来。若空间区域具有非平凡拓扑，全局单一光滑矢势可能不存在；本节只需使用上述局部关系。

\notetopic{QM-C02-S02-T03}{从 Lagrangian 完整恢复 Lorentz force}

令 $x_i$ 为 Cartesian coordinates、$v_i=\dot x_i$，并把式~\eqref{eq:QM-C02-S02-em-lagrangian} 写成
\[
  \mathcal L_{\mathrm{em}}
  =\frac12m\sum_jv_j^2-q\phi
  +\frac qc\sum_jv_jA_j.
\]
首先
\begin{equation}
  \frac{\partial\mathcal L_{\mathrm{em}}}{\partial v_i}
  =mv_i+\frac qcA_i.
  \label{eq:QM-C02-S02-dl-dv}
\end{equation}
由于 $A_i=A_i(\bm r(t),t)$，沿粒子轨迹的全导数必须同时包含显式时间变化与位置变化：
\begin{equation}
  \frac{dA_i}{dt}
  =\frac{\partial A_i}{\partial t}
  +\sum_jv_j\frac{\partial A_i}{\partial x_j}.
  \label{eq:QM-C02-S02-total-derivative-A}
\end{equation}
因此
\begin{equation}
  \frac{d}{dt}
  \left(\frac{\partial\mathcal L_{\mathrm{em}}}{\partial v_i}\right)
  =m\dot v_i+\frac qc\frac{\partial A_i}{\partial t}
  +\frac qc\sum_jv_j\frac{\partial A_i}{\partial x_j}.
  \label{eq:QM-C02-S02-dt-dl-dv}
\end{equation}
另一方面，求 $\partial/\partial x_i$ 时各 $v_j$ 作为独立变量保持不变，故
\begin{equation}
  \frac{\partial\mathcal L_{\mathrm{em}}}{\partial x_i}
  =-q\frac{\partial\phi}{\partial x_i}
  +\frac qc\sum_jv_j\frac{\partial A_j}{\partial x_i}.
  \label{eq:QM-C02-S02-dl-dx}
\end{equation}
代入 Euler--Lagrange 方程并整理：
\begin{align}
  m\dot v_i
  &=-q\frac{\partial\phi}{\partial x_i}
    -\frac qc\frac{\partial A_i}{\partial t}
    +\frac qc\sum_jv_j
      \left(
        \frac{\partial A_j}{\partial x_i}
        -\frac{\partial A_i}{\partial x_j}
      \right)\notag\\
  &=qE_i+\frac qc(\bm v\times\bm B)_i.
  \label{eq:QM-C02-S02-lorentz-components}
\end{align}
最后一步使用
\begin{equation}
  [\bm v\times(\nabla\times\bm A)]_i
  =\sum_jv_j
  \left(
    \frac{\partial A_j}{\partial x_i}
    -\frac{\partial A_i}{\partial x_j}
  \right).
  \label{eq:QM-C02-S02-curl-identity}
\end{equation}
三个分量合并后正好恢复式~\eqref{eq:QM-C02-S02-lorentz-force}。推导中最容易遗漏的是式~\eqref{eq:QM-C02-S02-total-derivative-A} 的对流项 $\sum_jv_j\partial A_i/\partial x_j$；没有它就无法形成 curl 的反对称导数组合。

\notetopic{QM-C02-S02-T04}{正则动量、机械动量与广义势}

由式~\eqref{eq:QM-C02-S02-dl-dv}，canonical momentum 为
\begin{equation}
  \boxed{
  \bm p=\frac{\partial\mathcal L_{\mathrm{em}}}{\partial\bm v}
  =m\bm v+\frac qc\bm A.}
  \label{eq:QM-C02-S02-canonical-momentum}
\end{equation}
mechanical momentum 记作 $\bm\pi=m\bm v$，所以
\begin{equation}
  \boxed{\bm\pi=\bm p-\frac qc\bm A.}
  \label{eq:QM-C02-S02-mechanical-momentum}
\end{equation}
规范变换会使 $\bm p'=\bm p+q\nabla\chi/c$，但式~\eqref{eq:QM-C02-S02-mechanical-momentum} 中的组合保持不变。正则动量是与坐标共轭的变量，机械动量则直接由速度决定；在电磁场中不能把符号 $\bm p$ 自动替换成 $m\bm v$。

形式上可定义 velocity-dependent generalized potential（速度依赖广义势）
\begin{equation}
  U(\bm r,\bm v,t)
  =q\phi-\frac qc\bm v\cdot\bm A,
  \qquad
  \mathcal L_{\mathrm{em}}=T-U.
  \label{eq:QM-C02-S02-generalized-potential}
\end{equation}
对于这类 $U$，物理广义力不再是 $-\partial U/\partial q_i$，而是
\begin{equation}
  \boxed{
  Q_i^{\mathrm{appl}}
  =\frac{d}{dt}\left(\frac{\partial U}{\partial\dot q_i}\right)
  -\frac{\partial U}{\partial q_i}.}
  \label{eq:QM-C02-S02-velocity-potential-force}
\end{equation}
磁力恒满足 $(q\bm v\times\bm B/c)\cdot\bm v=0$，会改变轨迹方向但不做功。因此 $-q\bm v\cdot\bm A/c$ 不能解释为普通磁势能；Lagrangian 本身也不是系统能量。

\notetopic{QM-C02-S02-T05}{相互作用、坐标几何、约束与非保守力}

Euler--Lagrange 方程右侧为零，不表示系统没有外力；它表示当前模型中由作用量描述的相互作用已经写入 $\mathcal L$。是否由外部给定，不是判断某个作用是否需要放在方程右侧的标准。例如给定电磁场虽可视为外场，其作用已经包含在式~\eqref{eq:QM-C02-S02-em-lagrangian} 中。

需要区分以下来源：
\begin{itemize}[leftmargin=*]
  \item \textbf{坐标几何：}曲线坐标使 $T$ 显含 $q_i$，产生的项用于组装该坐标系中的完整加速度，不是额外物理力。
  \item \textbf{速度依赖相互作用：}即使使用 Cartesian coordinates，$\bm v\cdot\bm A$ 仍会把正则动量从 $m\bm v$ 移开；这是相互作用本身造成的。
  \item \textbf{理想完整约束：}可用独立广义坐标消去约束，或用 Lagrange multiplier（拉格朗日乘子）保留并求出约束力。
  \item \textbf{未写入 $\mathcal L$ 的非保守作用：}使用 Lagrange--d'Alembert equation
  \begin{equation}
    \frac{d}{dt}\left(\frac{\partial\mathcal L}{\partial\dot q_i}\right)
    -\frac{\partial\mathcal L}{\partial q_i}
    =Q_i^{\mathrm{nc}}.
    \label{eq:QM-C02-S02-lagrange-dalembert}
  \end{equation}
\end{itemize}
仅写 $V=0$ 不能推出没有物理力，还必须排除未建模的外加作用与约束力。反过来，非零角动量也不意味着存在力：不经过坐标原点的自由直线运动可以具有非零且守恒的角动量。固定长度绳子的建模见\relatedexercise{QM-C02-S02-E02}{自编检验题 2}。

\subsection{适用条件与常见误区}

\begin{itemize}[leftmargin=*]
  \item 教材本节使用 Gaussian units；与 SI 公式比较时必须同步处理 $1/c$。
  \item $dA_i/dt$ 是沿粒子轨迹的全导数，不能误写成 $\partial A_i/\partial t$。
  \item $\bm B=\nabla\times\bm A$ 与式~\eqref{eq:QM-C02-S02-scalar-potential} 来自齐次 Maxwell 方程，不是为匹配 Lorentz force 临时定义的。
  \item 正则动量依赖规范，机械动量及由 $\bm E,\bm B$ 决定的运动不依赖规范。
  \item $V=0$ 只排除了所写的普通势能，不能自动排除约束力、摩擦力或其他未建模作用。
\end{itemize}

\subsection{一分钟回顾}

齐次 Maxwell 方程允许用 $\phi$ 与 $\bm A$ 表示电磁场；电磁 Lagrangian 中的 $\bm v\cdot\bm A$ 项通过反对称空间导数组合产生磁力。正则动量为 $\bm p=m\bm v+q\bm A/c$，而可直接由速度确定的机械动量为 $m\bm v$。规范变换只让 Lagrangian 增加全时间导数，不改变场和运动。坐标几何、速度依赖相互作用、理想约束与未建模非保守力必须在建模时分别处理。

\section{第 2.3 节：二体问题}
\label{section:QM-C02-S03}

\studymeta
  {QM-C02-S03}
  {2026-08-29}
  {Shankar, \textit{Principles of Quantum Mechanics}, Second Edition}
  {第 2 章 \S 2.3；印刷页 pp.~85--86；PDF pp.~98--99}
  {正文及教材练习 2.3.1 已核对；坐标变换、约化质量与适用条件已补全}

\subsection{本节主线}

两个只通过内部势相互作用、且不受外力的粒子，表面上需要同时求解 $\bm r_1(t)$ 与 $\bm r_2(t)$。改用 center-of-mass coordinate（质心坐标）$\bm R$ 和 relative coordinate（相对坐标）$\bm r$ 后，Lagrangian 可以分成互不耦合的质心运动与相对运动。原二体问题由此化为两个单粒子问题：一个质量为总质量 $M$ 的自由质心粒子，以及一个质量为 reduced mass（约化质量）$\mu$、在势 $V(\bm r)$ 中运动的等效粒子。

\notetopic{QM-C02-S03-T01}{建模假设与质心--相对坐标}

设两粒子的质量和位置分别为 $m_1,m_2$ 与 $\bm r_1,\bm r_2$。本节假设没有外力，且内部势只依赖两粒子的相对位置：
\begin{equation}
  V=V(\bm r_1-\bm r_2).
  \label{eq:QM-C02-S03-relative-potential}
\end{equation}
定义总质量、相对坐标和质心坐标
\begin{equation}
  M=m_1+m_2,
  \qquad
  \bm r=\bm r_1-\bm r_2,
  \qquad
  \bm R=\frac{m_1\bm r_1+m_2\bm r_2}{M}.
  \label{eq:QM-C02-S03-coordinate-definitions}
\end{equation}
解这个线性方程组可得逆变换
\begin{equation}
  \boxed{
  \bm r_1=\bm R+\frac{m_2}{M}\bm r,
  \qquad
  \bm r_2=\bm R-\frac{m_1}{M}\bm r.}
  \label{eq:QM-C02-S03-inverse-transform}
\end{equation}
例如 $m_1=3m,m_2=m$ 时，$\bm r_1=\bm R+\bm r/4$、$\bm r_2=\bm R-3\bm r/4$；较重的粒子离质心较近。

若把整个系统平移同一向量 $\bm a$，则
\begin{equation}
  \bm r_1\mapsto\bm r_1+\bm a,
  \quad
  \bm r_2\mapsto\bm r_2+\bm a,
  \quad
  \bm R\mapsto\bm R+\bm a,
  \quad
  \bm r\mapsto\bm r.
  \label{eq:QM-C02-S03-translation}
\end{equation}
因此 $V(\bm r)$ 对整体平移不敏感。能够分离质心运动的关键不是“势能不含 $\dot{\bm r}$”，而是势能不含 $\bm R$；前者对普通位置势本来就成立，后者才表达系统的平移不变性。

\notetopic{QM-C02-S03-T02}{动能分离与约化质量}

原 Lagrangian 为
\begin{equation}
  \mathcal L
  =\frac12m_1|\dot{\bm r}_1|^2
  +\frac12m_2|\dot{\bm r}_2|^2
  -V(\bm r_1-\bm r_2).
  \label{eq:QM-C02-S03-original-lagrangian}
\end{equation}
对式~\eqref{eq:QM-C02-S03-inverse-transform} 求时间导数：
\begin{equation}
  \dot{\bm r}_1=\dot{\bm R}+\frac{m_2}{M}\dot{\bm r},
  \qquad
  \dot{\bm r}_2=\dot{\bm R}-\frac{m_1}{M}\dot{\bm r}.
  \label{eq:QM-C02-S03-velocity-transform}
\end{equation}
代入动能并展开，第一部分为
\begin{align}
  \frac12m_1|\dot{\bm r}_1|^2
  ={}&\frac12m_1|\dot{\bm R}|^2
  +\frac{m_1m_2}{M}\dot{\bm R}\cdot\dot{\bm r}
  +\frac12m_1\frac{m_2^2}{M^2}|\dot{\bm r}|^2,
  \label{eq:QM-C02-S03-first-kinetic-expansion}
\end{align}
第二部分为
\begin{align}
  \frac12m_2|\dot{\bm r}_2|^2
  ={}&\frac12m_2|\dot{\bm R}|^2
  -\frac{m_1m_2}{M}\dot{\bm R}\cdot\dot{\bm r}
  +\frac12m_2\frac{m_1^2}{M^2}|\dot{\bm r}|^2.
  \label{eq:QM-C02-S03-second-kinetic-expansion}
\end{align}
交叉项恰好抵消。相对速度平方项的系数满足
\begin{align}
  m_1\frac{m_2^2}{M^2}+m_2\frac{m_1^2}{M^2}
  &=\frac{m_1m_2(m_1+m_2)}{M^2}
   =\frac{m_1m_2}{m_1+m_2}.
  \label{eq:QM-C02-S03-reduced-mass-coefficient}
\end{align}
定义
\begin{equation}
  \boxed{
  \mu=\frac{m_1m_2}{m_1+m_2},
  \qquad
  \frac1\mu=\frac1{m_1}+\frac1{m_2}.}
  \label{eq:QM-C02-S03-reduced-mass}
\end{equation}
于是
\begin{equation}
  \boxed{
  \mathcal L
  =\frac12M|\dot{\bm R}|^2
  +\frac12\mu|\dot{\bm r}|^2
  -V(\bm r).}
  \label{eq:QM-C02-S03-separated-lagrangian}
\end{equation}
教材练习 2.3.1 要求完成这一变量代换，完整计算见\relatedexercise{QM-C02-S03-E01}{教材练习 2.3.1}。

\notetopic{QM-C02-S03-T03}{两个互不耦合的单粒子问题}

式~\eqref{eq:QM-C02-S03-separated-lagrangian} 可以写成
\begin{equation}
  \mathcal L
  =\underbrace{\frac12M|\dot{\bm R}|^2}_{\mathcal L_{\mathrm{CM}}}
  +\underbrace{\left(\frac12\mu|\dot{\bm r}|^2-V(\bm r)\right)}_{\mathcal L_{\mathrm{rel}}}.
  \label{eq:QM-C02-S03-lagrangian-parts}
\end{equation}
质心坐标不显含于 $\mathcal L$，所以它是循环坐标，其共轭动量为
\begin{equation}
  \boxed{
  \bm P=\frac{\partial\mathcal L}{\partial\dot{\bm R}}
  =M\dot{\bm R}
  =m_1\dot{\bm r}_1+m_2\dot{\bm r}_2,}
  \qquad
  \dot{\bm P}=0.
  \label{eq:QM-C02-S03-total-momentum}
\end{equation}
这正是系统总动量。相对坐标的共轭动量和两组运动方程为
\begin{equation}
  \boxed{
  \bm p=\mu\dot{\bm r},
  \qquad
  M\ddot{\bm R}=0,
  \qquad
  \mu\ddot{\bm r}=-\nabla_{\bm r}V(\bm r).}
  \label{eq:QM-C02-S03-equations-of-motion}
\end{equation}
因此第一个等效粒子质量为 $M$，做自由质心运动；第二个等效粒子质量为 $\mu$，在势 $V(\bm r)$ 中运动。它们只是对原坐标的等效描述，并不表示系统中真的新增了两个粒子。

在 center-of-mass frame（质心参考系）中取 $\bm R=0$、$\dot{\bm R}=0$，则
\begin{equation}
  \dot{\bm r}_1=\frac{m_2}{M}\dot{\bm r},
  \qquad
  \dot{\bm r}_2=-\frac{m_1}{M}\dot{\bm r},
  \label{eq:QM-C02-S03-cm-velocities}
\end{equation}
并且
\begin{equation}
  m_1\dot{\bm r}_1=\mu\dot{\bm r}=\bm p,
  \qquad
  m_2\dot{\bm r}_2=-\mu\dot{\bm r}=-\bm p.
  \label{eq:QM-C02-S03-cm-momenta}
\end{equation}
研究内部运动时可以忽略平凡的质心匀速运动，但这不表示质心运动不存在；只是在质心参考系中把它消去了。

\notetopic{QM-C02-S03-T04}{从 Newton 方程理解约化质量}

设粒子 1 受到的内部力为 $\bm F$，粒子 2 受到的内部力为 $-\bm F$。由
\begin{equation}
  m_1\ddot{\bm r}_1=\bm F,
  \qquad
  m_2\ddot{\bm r}_2=-\bm F
  \label{eq:QM-C02-S03-newton-pair}
\end{equation}
可得
\begin{align}
  \ddot{\bm r}
  &=\ddot{\bm r}_1-\ddot{\bm r}_2
    =\left(\frac1{m_1}+\frac1{m_2}\right)\bm F
    =\frac{\bm F}{\mu},
\end{align}
即
\begin{equation}
  \boxed{\mu\ddot{\bm r}=\bm F.}
  \label{eq:QM-C02-S03-newton-reduced-mass}
\end{equation}
若 $\bm F=-\nabla_{\bm r}V$，这正是式~\eqref{eq:QM-C02-S03-equations-of-motion} 的相对运动方程。约化质量把两个物体共同改变相对距离的效果压缩进一个等效惯性参数中。因为
\[
  \frac1\mu=\frac1{m_1}+\frac1{m_2},
\]
所以 $\mu<\min(m_1,m_2)$。两个常用极限是
\begin{equation}
  m_2\gg m_1\Longrightarrow\mu\simeq m_1,
  \qquad
  m_1=m_2=m\Longrightarrow\mu=\frac m2.
  \label{eq:QM-C02-S03-reduced-mass-limits}
\end{equation}
第一种情形表示重粒子近似固定，轻粒子近似以自身质量在势中运动；第二种情形中两粒子的加速度大小相等、方向相反，因此相对加速度是任一粒子加速度的两倍。$\mu=m/2$ 仍只是相对运动的等效质量，不是新的真实粒子。

\notetopic{QM-C02-S03-T05}{分离成立的条件}

坐标变换式~\eqref{eq:QM-C02-S03-coordinate-definitions} 总可以进行，但动力学分离还依赖模型结构：
\begin{itemize}[leftmargin=*]
  \item 内部相互作用只依赖 $\bm r_1-\bm r_2$，因而不显含 $\bm R$；
  \item 系统没有把 $\bm R$ 与 $\bm r$ 重新耦合的一般外势；
  \item 质心与相对坐标的动能交叉项因质量加权的质心定义而抵消。
\end{itemize}
例如，只作用于粒子 1 的外势会变为
\begin{equation}
  U(\bm r_1)
  =U\left(\bm R+\frac{m_2}{M}\bm r\right),
  \label{eq:QM-C02-S03-external-potential}
\end{equation}
一般同时依赖 $\bm R$ 与 $\bm r$，从而破坏分离。这里必须保留“一般”二字：特殊外势仍可能分解成纯质心项与纯相对项。另外，二体分离只要求 $V=V(\bm r)$；更强的 central potential 条件 $V=V(|\bm r|)$ 还会进一步带来相对运动的旋转对称性和角动量守恒。

\subsection{适用条件与常见误区}

\begin{itemize}[leftmargin=*]
  \item 约化质量是相对坐标的等效惯性参数，不是任何一个真实粒子的质量。
  \item $M\ddot{\bm R}=0$ 表示质心做匀速运动，不表示质心坐标可以在所有参考系中直接删除。
  \item 势能不含 $\dot{\bm r}$ 不是分离的理由；关键是 $V$ 不含 $\bm R$。
  \item 改变量本身不减少自由度，而是把原先耦合的自由度改写成更合适的独立组合。
  \item 一般外势会耦合 $\bm R$ 与 $\bm r$；是否仍能分离必须代入具体势后判断。
\end{itemize}

\subsection{一分钟回顾}

质量加权的质心坐标与相对坐标把二体 Lagrangian 对角化为质心部分和相对部分。质心以总质量 $M$ 做自由匀速运动，总动量守恒；相对坐标则等价于一个质量为 $\mu=m_1m_2/(m_1+m_2)$ 的粒子在势 $V(\bm r)$ 中运动。约化质量既可从动能展开得到，也可从两条 Newton 方程相减得到。真正保证分离的是平移不变的内部势与无一般外势，而不是简单地把二体问题“少算一个粒子”。

\section{第 2.4 节：粒子有多聪明？}
\label{section:QM-C02-S04}

\studymeta
  {QM-C02-S04}
  {2026-08-30}
  {Shankar, \textit{Principles of Quantum Mechanics}, Second Edition}
  {第 2 章 \S 2.4；印刷页 p.~86；PDF p.~99}
  {正文已核对；本节无教材例题或习题，配套题明确标为自编检验题}

\subsection{本节主线}

作用量以整条路径为输入，固定起点与终点后再要求 $\delta S=0$，形式上似乎暗示粒子必须预先知道终点、比较所有路径后才选择运动方式。本节指出这种 teleological（目的论式）理解是错觉：全局的变分表述会导出局部 Euler--Lagrange 方程，经典粒子只需服从当前时刻的局部运动规律。教材随后预告 path integral（路径积分）：量子传播对所有路径的复振幅求和，而经典路径在适当极限中由 stationary phase（平稳相位）机制显现。

\notetopic{QM-C02-S04-T01}{全局变分不等于粒子预知未来}

对固定端点 $x(t_i)=x_i$、$x(t_f)=x_f$，作用量为
\begin{equation}
  S[x]=\int_{t_i}^{t_f}
  \mathcal L(x,\dot x,t)\,dt.
  \label{eq:QM-C02-S04-action}
\end{equation}
允许变分满足
\[
  x(t)\mapsto x(t)+\epsilon\eta(t),
  \qquad
  \eta(t_i)=\eta(t_f)=0.
\]
第 2.1 节的推导给出
\begin{equation}
  \delta S
  =\int_{t_i}^{t_f}
  \left[
    \frac{\partial\mathcal L}{\partial x}
    -\frac{d}{dt}
      \left(\frac{\partial\mathcal L}{\partial\dot x}\right)
  \right]\eta(t)\,dt.
  \label{eq:QM-C02-S04-first-variation}
\end{equation}
由于该式对所有允许的 $\eta(t)$ 都为零，得到局部方程
\begin{equation}
  \boxed{
  \frac{d}{dt}
  \left(\frac{\partial\mathcal L}{\partial\dot x}\right)
  -\frac{\partial\mathcal L}{\partial x}=0.}
  \label{eq:QM-C02-S04-euler-lagrange}
\end{equation}
因此，“对整条路径取变分”是描述运动规律的数学形式，不表示粒子在物理上逐条计算候选路径。

需要区分两种求解方式：
\begin{itemize}[leftmargin=*]
  \item \textbf{Initial-value problem（初值问题）：}给定 $x(t_i)$ 与 $\dot x(t_i)$，用局部微分方程向未来演化；
  \item \textbf{Boundary-value problem（边值问题）：}给定 $x(t_i)$ 与 $x(t_f)$，寻找满足这些边界条件的平稳路径。
\end{itemize}
未来终点是求解者对边值问题施加的条件，用于从可能的运动中选出符合边界的轨迹；它不是粒子获得的未来信息。对
\begin{equation}
  \mathcal L=\frac12m\dot x^2-V(x,t)
  \label{eq:QM-C02-S04-standard-lagrangian}
\end{equation}
有
\begin{equation}
  \boxed{
  m\ddot x=-\frac{\partial V(x,t)}{\partial x}.}
  \label{eq:QM-C02-S04-local-newton}
\end{equation}
当前加速度只由当前的 $x,t$ 和该处势能梯度决定；在更一般的速度依赖 Lagrangian 中，局部运动还可能依赖当前速度。

\notetopic{QM-C02-S04-T02}{Stationary 不等于 minimum}

教材沿用 principle of least action（最小作用量原理）的传统名称，并在本节使用“选择最小作用量路径”的直观措辞。严格条件仍然是
\begin{equation}
  \boxed{\delta S=0,}
  \label{eq:QM-C02-S04-stationary-action}
\end{equation}
即作用量在经典路径处 stationary（平稳）。它可能是局部最小值、局部最大值或 saddle point（鞍点）；仅有一阶变分为零不能判断二阶性质。因此本节的 least/minimize 应理解为传统而宽松的表述。

教材还说粒子只需探测附近的势并沿变化最快的方向加速。对 Cartesian coordinates 下的普通保守系统
\[
  \mathcal L=\frac12m|\dot{\bm r}|^2-V(\bm r,t)
\]
确有 $m\ddot{\bm r}=-\nabla V$。但电磁场、曲线坐标、约束或速度依赖相互作用中的局部 Euler--Lagrange 方程不一定能简化为这一形式。

\notetopic{QM-C02-S04-T03}{路径积分预告与经典路径的显现}

本节只作概念预告；后续 Feynman path integral 的标准形式为
\begin{equation}
  K(x_f,t_f;x_i,t_i)
  =\int\mathcal D x(t)\,
  \exp\left(\frac{i}{\hbar}S[x]\right).
  \label{eq:QM-C02-S04-path-integral}
\end{equation}
$K$ 是从初始时空点传播到末时空点的 quantum amplitude（量子振幅），$\int\mathcal D x(t)$ 形式上表示对所有连接端点的路径求和。对普通实作用量，单条路径的因子满足
\begin{equation}
  \left|
    \exp\left(\frac{i}{\hbar}S[x]\right)
  \right|=1,
  \label{eq:QM-C02-S04-equal-modulus}
\end{equation}
所以教材所说的 equal weight 应理解为相同模长的复振幅，而不是各路径具有相同经典概率；不同路径的相位 $S[x]/\hbar$ 仍然不同。

在经典路径 $x_{\mathrm{cl}}$ 附近令 $x=x_{\mathrm{cl}}+\eta$，有泛函 Taylor 展开
\begin{equation}
  S[x_{\mathrm{cl}}+\eta]
  =S[x_{\mathrm{cl}}]
  +\delta S
  +\frac12\delta^2S+\cdots.
  \label{eq:QM-C02-S04-action-expansion}
\end{equation}
由于 $\delta S[x_{\mathrm{cl}}]=0$，经典路径附近的作用量差从 $O(\eta^2)$ 开始，相位变化较慢，邻近路径的振幅较容易相干叠加。非平稳路径附近一般存在 $O(\eta)$ 的作用量变化；当典型作用量满足
\begin{equation}
  \frac{S}{\hbar}\gg1
  \label{eq:QM-C02-S04-classical-limit}
\end{equation}
时，邻近路径的相位会快速改变，其贡献大多相消。因而经典极限中起主要作用的不是一条孤立路径，而是平稳作用量路径附近的一簇相干贡献。

“粒子同时沿所有路径运动”是对路径积分计算形式的启发式说法。路径积分首先是一种计算量子传播振幅的形式体系，不能未经进一步解释就把每条积分路径视为一次可观测的经典轨迹。相关概念检查见\relatedexercise{QM-C02-S04-E01}{自编检验题 1}。

\subsection{适用条件与常见误区}

\begin{itemize}[leftmargin=*]
  \item 固定未来终点是边值问题的数学条件，不代表粒子拥有未来信息。
  \item $\delta S=0$ 是平稳条件，不能自动加强为全局最小值。
  \item equal weight 不等于 equal probability；路径积分相加的是带不同相位的复振幅。
  \item 非平稳路径的强烈相消依赖经典极限 $S/\hbar\gg1$；量子尺度下不能直接忽略多路径干涉。
  \item 式~\eqref{eq:QM-C02-S04-path-integral} 在此只是后续知识的方向性预告，路径测度等严格问题留待正式学习路径积分时处理。
\end{itemize}

\subsection{一分钟回顾}

平稳作用量以整条路径表述问题，却导出只依赖局部状态的 Euler--Lagrange 方程，因此不要求经典粒子预知未来。严格条件是 stationary action，而不一定是 minimum action。量子路径积分对所有路径的复振幅求和；不同相位使非平稳路径在经典极限中大多相消，而平稳作用量路径附近的贡献能够相干叠加，从而显现经典轨迹。

\section{第 2.5 节：Hamilton 形式}
\label{section:QM-C02-S05}

\studymeta
  {QM-C02-S05}
  {2026-08-31}
  {Shankar, \textit{Principles of Quantum Mechanics}, Second Edition}
  {第 2 章 \S 2.5；印刷页 pp.~86--90；PDF pp.~99--103}
  {正文、教材讲解例题、练习 2.5.1--2.5.4 与收口检查已核对}

\subsection{本节主线}

Lagrange 形式把广义坐标 $q_i$ 与广义速度 $\dot q_i$ 作为基本变量，再由
\[
  p_i=\frac{\partial\mathcal L}{\partial\dot q_i}
\]
导出正则动量。Hamilton 形式对速度变量作 Legendre transformation（Legendre 变换），改用 $(q_i,p_i)$ 描述同一动力学：$n$ 条二阶 Euler--Lagrange 方程被改写为 $2n$ 条一阶 Hamilton 正则方程。系统状态因而成为 $2n$ 维 phase space（相空间）中的一点，而 Hamiltonian 决定该点如何流动。

需要分开判断两个命题：Hamiltonian 是否等于通常的机械能，以及 Hamiltonian 是否守恒。前者依赖 Lagrangian 对速度的结构，后者依赖是否显含时间；两者不能混为一谈。

\notetopic{QM-C02-S05-T01}{从构型空间到相空间}

对 $n$ 个广义坐标，Lagrange 形式使用
\begin{equation}
  \mathcal L=\mathcal L(q_1,\ldots,q_n,
  \dot q_1,\ldots,\dot q_n,t),
  \label{eq:QM-C02-S05-lagrangian-variables}
\end{equation}
其中 $q_i,\dot q_i$ 是独立变量，正则动量由
\begin{equation}
  p_i=\frac{\partial\mathcal L}{\partial\dot q_i}
  \label{eq:QM-C02-S05-canonical-momentum}
\end{equation}
定义。Hamilton 形式则使用
\begin{equation}
  \mathcal H=\mathcal H(q_1,\ldots,q_n,p_1,\ldots,p_n,t),
  \label{eq:QM-C02-S05-hamiltonian-variables}
\end{equation}
并由运动方程求出 $\dot q_i,\dot p_i$。

只给构型空间中的位置 $(q_1,\ldots,q_n)$，一般不能确定下一时刻的运动，因为同一位置可以具有不同速度。相空间坐标
\begin{equation}
  (q_1,\ldots,q_n,p_1,\ldots,p_n)
  \label{eq:QM-C02-S05-phase-point}
\end{equation}
在 Legendre 变换可逆时等价于 $(q_i,\dot q_i)$，给出了完整的瞬时状态。$n$ 条二阶方程需要 $q_i(t_0),\dot q_i(t_0)$ 共 $2n$ 个初始数据；$2n$ 条一阶方程需要 $q_i(t_0),p_i(t_0)$，初始数据总数不变。

正则动量是式~\eqref{eq:QM-C02-S05-canonical-momentum} 定义的共轭变量，不应预设为 $m_i\dot q_i$。它只在普通 Cartesian kinetic energy 等特定情形下等于相应机械动量；速度依赖相互作用会改变二者关系。

\notetopic{QM-C02-S05-T02}{Legendre 变换及其可逆条件}

先考虑单变量函数 $f(x)$。定义新变量与新函数
\begin{equation}
  u=f'(x),
  \qquad
  g(u)=x(u)u-f(x(u)).
  \label{eq:QM-C02-S05-legendre-one-variable}
\end{equation}
若 $u=f'(x)$ 可以局部反解为 $x=x(u)$，则
\begin{align}
  \frac{\dd g}{\dd u}
  &=x+u\frac{\dd x}{\dd u}
    -f'(x)\frac{\dd x}{\dd u}
    =x.
  \label{eq:QM-C02-S05-legendre-derivative}
\end{align}
因为 $f'(x)=u$，含 $\dd x/\dd u$ 的两项正好抵消。Legendre 变换以切线斜率 $u$ 代替切点坐标 $x$ 描述同一函数；$g(u)$ 等于该切线纵截距的相反数。

例如，对
\[
  f(x)=\frac12ax^2,
  \qquad a\neq0,
\]
有
\begin{equation}
  u=ax,
  \qquad
  x=\frac ua,
  \qquad
  g(u)=\frac{u^2}{2a}.
  \label{eq:QM-C02-S05-quadratic-legendre}
\end{equation}

Hamiltonian 是只对所有速度变量进行的多变量 Legendre 变换：
\begin{equation}
  \boxed{
  \mathcal H(q,p,t)
  =\sum_{i=1}^{n}p_i\dot q_i-\mathcal L(q,\dot q,t),
  \qquad
  p_i=\frac{\partial\mathcal L}{\partial\dot q_i}.}
  \label{eq:QM-C02-S05-hamiltonian-definition}
\end{equation}
定义完成后，必须把每个 $\dot q_i$ 消去，写成 $\dot q_i(q,p,t)$；否则所得表达式还不是只以 $(q,p,t)$ 为独立变量的普通 Hamiltonian。

\dialoguesupplement{教材没有在此展开变换的正则性条件。严格地说，需要速度 Hessian matrix
\begin{equation}
  W_{ij}=\frac{\partial^2\mathcal L}
  {\partial\dot q_i\partial\dot q_j}
  =\frac{\partial p_i}{\partial\dot q_j}
  \label{eq:QM-C02-S05-velocity-hessian}
\end{equation}
在所研究区域可逆。局部充分条件是 $\det W\neq0$。若 $\det W=0$，映射 $\dot q\mapsto p$ 可能无法反解，系统需要带约束的 Hamilton 形式，不能直接套用本节普通推导。}

例如
\begin{equation}
  \mathcal L=a(q)\dot q-V(q)
  \label{eq:QM-C02-S05-singular-lagrangian}
\end{equation}
给出
\[
  p=a(q),
  \qquad
  \frac{\partial p}{\partial\dot q}=0.
\]
$p$ 完全不含 $\dot q$，所有速度被映到同一个动量值，因而不能由 $(q,p)$ 唯一恢复 $\dot q$。

\notetopic{QM-C02-S05-T03}{Hamilton 正则方程的完整推导}

由 $\mathcal L=\mathcal L(q,\dot q,t)$，其全微分为
\begin{equation}
  \dd\mathcal L
  =\sum_i\frac{\partial\mathcal L}{\partial q_i}\dd q_i
  +\sum_i p_i\dd\dot q_i
  +\frac{\partial\mathcal L}{\partial t}\dd t.
  \label{eq:QM-C02-S05-differential-lagrangian}
\end{equation}
对式~\eqref{eq:QM-C02-S05-hamiltonian-definition} 求全微分：
\begin{align}
  \dd\mathcal H
  &=\sum_i\dot q_i\dd p_i
    +\sum_i p_i\dd\dot q_i-\dd\mathcal L \notag\\
  &=\sum_i\dot q_i\dd p_i
    -\sum_i\frac{\partial\mathcal L}{\partial q_i}\dd q_i
    -\frac{\partial\mathcal L}{\partial t}\dd t.
  \label{eq:QM-C02-S05-differential-hamiltonian-first}
\end{align}
第二行中两组 $\sum_i p_i\dd\dot q_i$ 完全抵消。这正是 Legendre 变换把速度换成动量后，不再留下 $\dd\dot q_i$ 项的原因。

Euler--Lagrange 方程与式~\eqref{eq:QM-C02-S05-canonical-momentum} 给出
\begin{equation}
  \frac{\partial\mathcal L}{\partial q_i}
  =\frac{\dd}{\dd t}
    \left(\frac{\partial\mathcal L}{\partial\dot q_i}\right)
  =\dot p_i.
  \label{eq:QM-C02-S05-el-momentum}
\end{equation}
因此
\begin{equation}
  \dd\mathcal H
  =\sum_i\dot q_i\dd p_i
  -\sum_i\dot p_i\dd q_i
  -\frac{\partial\mathcal L}{\partial t}\dd t.
  \label{eq:QM-C02-S05-differential-hamiltonian-second}
\end{equation}
另一方面，$\mathcal H=\mathcal H(q,p,t)$ 的全微分是
\begin{equation}
  \dd\mathcal H
  =\sum_i\frac{\partial\mathcal H}{\partial q_i}\dd q_i
  +\sum_i\frac{\partial\mathcal H}{\partial p_i}\dd p_i
  +\frac{\partial\mathcal H}{\partial t}\dd t.
  \label{eq:QM-C02-S05-differential-hamiltonian-general}
\end{equation}
逐项比较独立微分 $\dd q_i,\dd p_i,\dd t$ 的系数，得到
\begin{equation}
  \boxed{
  \dot q_i=\frac{\partial\mathcal H}{\partial p_i},
  \qquad
  \dot p_i=-\frac{\partial\mathcal H}{\partial q_i},}
  \label{eq:QM-C02-S05-hamilton-equations}
\end{equation}
以及
\begin{equation}
  \boxed{
  \frac{\partial\mathcal H}{\partial t}
  =-\frac{\partial\mathcal L}{\partial t}.}
  \label{eq:QM-C02-S05-explicit-time}
\end{equation}
第一条正则方程从动量恢复速度，第二条给出动量的演化。若 Legendre 变换正则，这 $2n$ 条一阶方程与原来的 $n$ 条二阶 Euler--Lagrange 方程等价。

若 $\mathcal H$ 不显含某个 $q_k$，则
\begin{equation}
  \dot p_k=-\frac{\partial\mathcal H}{\partial q_k}=0,
  \label{eq:QM-C02-S05-cyclic-hamiltonian}
\end{equation}
所以对应正则动量守恒。这是循环坐标在 Hamilton 形式中的直接表达。

\notetopic{QM-C02-S05-T04}{Hamiltonian、能量与守恒必须分开判断}

考虑 standard mechanical Lagrangian
\begin{equation}
  \mathcal L(q,\dot q,t)=T(q,\dot q)-V(q,t),
  \label{eq:QM-C02-S05-standard-mechanical-lagrangian}
\end{equation}
其中 $V$ 不依赖速度，并且 $T$ 关于所有 $\dot q_i$ 是二次齐次函数，例如
\begin{equation}
  T=\frac12\sum_{i,j}M_{ij}(q)\dot q_i\dot q_j.
  \label{eq:QM-C02-S05-quadratic-kinetic-energy}
\end{equation}
由 Euler 齐次函数定理，
\begin{equation}
  \sum_i p_i\dot q_i
  =\sum_i\frac{\partial T}{\partial\dot q_i}\dot q_i
  =2T.
  \label{eq:QM-C02-S05-euler-homogeneous}
\end{equation}
所以
\begin{equation}
  \boxed{
  \mathcal H=2T-(T-V)=T+V.}
  \label{eq:QM-C02-S05-hamiltonian-energy}
\end{equation}
$M_{ij}$ 可以依赖坐标，因此该结论也适用于普通极坐标、球坐标等曲线坐标中的标准动能。若 Lagrangian 含一般的速度依赖相互作用，不能未经计算就断言 $\mathcal H=T+V$；必须从定义式~\eqref{eq:QM-C02-S05-hamiltonian-definition} 重新计算。

沿实际轨迹对 $\mathcal H$ 求时间导数，并使用正则方程：
\begin{align}
  \frac{\dd\mathcal H}{\dd t}
  &=\sum_i\frac{\partial\mathcal H}{\partial q_i}\dot q_i
    +\sum_i\frac{\partial\mathcal H}{\partial p_i}\dot p_i
    +\frac{\partial\mathcal H}{\partial t} \notag\\
  &=\sum_i(-\dot p_i)\dot q_i
    +\sum_i\dot q_i\dot p_i
    +\frac{\partial\mathcal H}{\partial t}
   =\frac{\partial\mathcal H}{\partial t}
   =-\frac{\partial\mathcal L}{\partial t}.
  \label{eq:QM-C02-S05-hamiltonian-time-derivative}
\end{align}
因此 $\partial\mathcal L/\partial t=0$ 保证 Hamiltonian 守恒，但只有在另行确认 $\mathcal H$ 表示物理能量后，才能把它直接称为能量守恒。反过来，即使式~\eqref{eq:QM-C02-S05-hamiltonian-energy} 成立，若 $V(q,t)$ 显含时间，$\mathcal H=T+V$ 一般仍不守恒。

\notetopic{QM-C02-S05-T05}{简谐振子与相空间轨迹}

一维简谐振子的 Lagrangian 是
\begin{equation}
  \mathcal L=\frac12m\dot x^2-\frac12kx^2.
  \label{eq:QM-C02-S05-oscillator-lagrangian}
\end{equation}
由 $p=m\dot x$、$\dot x=p/m$ 得
\begin{equation}
  \boxed{
  \mathcal H(x,p)=\frac{p^2}{2m}+\frac12kx^2.}
  \label{eq:QM-C02-S05-oscillator-hamiltonian}
\end{equation}
Hamilton 正则方程为
\begin{equation}
  \dot x=\frac pm,
  \qquad
  \dot p=-kx.
  \label{eq:QM-C02-S05-oscillator-hamilton-equations}
\end{equation}
消去 $p$ 就恢复 $m\ddot x=-kx$。在相空间中，状态 $(x,p)$ 的瞬时流向是
\[
  (\dot x,\dot p)=\left(\frac pm,-kx\right).
\]
固定能量 $E$ 的轨迹满足
\begin{equation}
  \boxed{
  \frac{x^2}{2E/k}+\frac{p^2}{2mE}=1,}
  \label{eq:QM-C02-S05-oscillator-ellipse}
\end{equation}
是一条半轴平方为 $a^2=2E/k$、$b^2=2mE$ 的椭圆。椭圆上的不同点是相同能量下的不同运动状态，而不同椭圆对应不同能量。

\begin{center}
\begin{tikzpicture}[>=Stealth,scale=0.92]
  \draw[->] (-3.4,0) -- (3.6,0) node[right] {$x$};
  \draw[->] (0,-2.25) -- (0,2.35) node[above] {$p$};
  \draw[thick,noteBlue] (0,0) ellipse (2.7 and 1.55);
  \draw[->,thick,accentRed] (2.7,0.2) -- (2.7,-0.55);
  \draw[->,thick,accentRed] (-0.45,1.55) -- (0.55,1.55);
  \draw[->,thick,accentRed] (-2.7,-0.2) -- (-2.7,0.55);
  \draw[->,thick,accentRed] (0.45,-1.55) -- (-0.55,-1.55);
  \node[noteBlue,fill=white,inner sep=1pt] at (1.65,1.23)
    {$\mathcal H(x,p)=E$};
  \node[below right] at (2.7,0) {$(a,0)$};
  \node[left] at (0,1.55) {$(0,b)$};
\end{tikzpicture}

\textit{简谐振子的相空间能量椭圆。红色箭头由 Hamilton 正则方程决定；在右端点 $(a,0)$，$\dot p=-ka<0$，轨迹向下。}
\end{center}

完整推导见习题部分的\hyperref[exercise:QM-C02-S05-E00]{教材讲解例题}；练习 2.5.1--2.5.4 紧随其后。

\subsection{适用条件与常见误区}

\begin{itemize}[leftmargin=*]
  \item Hamiltonian 的定义是对速度变量的 Legendre 变换，不是把 $\mathcal L=T-V$ 中的负号机械改为正号。
  \item $\mathcal H$ 的最终表达式必须只含 $(q,p,t)$，不能遗留尚未消去的 $\dot q_i$。
  \item 普通 Hamilton 形式要求 $\dot q\mapsto p$ 可逆；速度 Hessian 奇异时不能直接套用本节推导。
  \item $\mathcal H=T+V$ 与 $\dd\mathcal H/\dd t=0$ 是两个不同命题，分别检查速度结构与显式时间依赖。
  \item 正则动量不一定等于机械动量；具体关系必须由 $p_i=\partial\mathcal L/\partial\dot q_i$ 计算。
  \item 二体系统总动量守恒只保证质心静止或匀速直线运动，不表示两个粒子之间没有相对运动。
\end{itemize}

\subsection{一分钟回顾}

Legendre 变换把广义速度替换为正则动量，使系统状态成为 $2n$ 维相空间中的点。Hamiltonian 的全微分与 Euler--Lagrange 方程共同给出 $\dot q_i=\partial\mathcal H/\partial p_i$、$\dot p_i=-\partial\mathcal H/\partial q_i$。对速度二次齐次、势能不依赖速度的标准系统，$\mathcal H=T+V$；若 Lagrangian 不显含时间，Hamiltonian 守恒。简谐振子在相空间中沿固定能量椭圆运动，显示 Hamiltonian 同时编码能量曲线与时间演化方向。
